\section{Feature Engineering}
\label{sec:features}

Every feature consumed by the model or reported in this document is declared
in \texttt{conf/features.yaml} and implemented in exactly one of four modules
(Table~\ref{tab:features}), generated directly from that file by
\texttt{scripts/export\_feature\_tables.py}.

\subsection{Point-in-time contracts}

Two distinct guarantees are needed, because the features fall into two
different classes.

\textbf{Corpus aggregates} (\texttt{item\_popularity\_count},
\texttt{visitor\_prior\_session\_count}) are fitted once on the training
partition and accept an explicit \texttt{cutoff\_order} argument so a strict
\texttt{<} comparison against a global event-order index can be asserted
directly: an event with \texttt{event\_order} $\geq$ the cut-off must not
change the value, and an event exactly \emph{at} the cut-off belongs to the
future, not the past.

\textbf{Session-local features}
(\texttt{session\_length\_so\_far}, \texttt{time\_since\_session\_start}) have
no corpus-level fit at all: they are computed per event row from only the
\emph{other} events in the same session that occurred strictly before it.
\texttt{time\_since\_session\_start} in particular reads the timestamp of the
most recently \emph{observed} event, not the timestamp of the row being
scored -- the latter is the quantity the model is trying to predict alongside
the next item, and using it would be the single most tempting leak in this
codebase.

\textbf{Context features} (\texttt{hour\_of\_day}, \texttt{is\_weekend}) are
deterministic functions of the event being scored and are declared
\texttt{point\_in\_time: false}: they carry no leakage risk because the
wall-clock time of a prediction request is known by definition, unlike the
item that will actually be clicked.

\subsection{Catalogue features}

\texttt{item\_category} is read directly from the item-properties changelog,
independent of any interaction history, so cold-start items -- never seen in
training -- still carry a category and are not excluded from candidate
scoring.

\subsection{An explicitly ablation-only visitor feature}

\texttt{visitor\_prior\_session\_count} counts how many training sessions a
\texttt{visitorid} started before the current one. The primary session-kNN
method never conditions on it: the whole premise of this project is that
anonymous, short-lived sessions can be modelled without a persistent
identity. The feature exists so that premise is testable (does a returning
cookie help at all?) instead of merely asserted; Section~\ref{sec:results}
reports the answer.

% Generated by scripts/export_feature_tables.py -- do not edit by hand.
\begin{table}[t]
  \centering
  \caption{Engineered features. Aggregates (\texttt{fit\_on: train\_only})
  are computed point-in-time against a strict event-order cut-off;
  session-local features (\texttt{fit\_on: none}) are point-in-time by
  construction, since they read only strictly-earlier positions in the
  same session.}
  \label{tab:features}
  \begin{tabular}{llll}
    \toprule
    Feature & Family & Fit scope & Point-in-time \\
    \midrule
    \texttt{item\_popularity\_count} & item & \texttt{train\_only} & \checkmark \\
    \texttt{item\_category} & item & \texttt{catalogue} & -- \\
    \texttt{visitor\_prior\_session\_count} & user & \texttt{train\_only} & \checkmark \\
    \texttt{session\_length\_so\_far} & interaction & \texttt{none} & \checkmark \\
    \texttt{time\_since\_session\_start} & interaction & \texttt{none} & \checkmark \\
    \texttt{hour\_of\_day} & context & \texttt{none} & -- \\
    \texttt{is\_weekend} & context & \texttt{none} & -- \\
    \bottomrule
  \end{tabular}
\end{table}


\subsection{Rationale}

% Generated by scripts/export_feature_tables.py -- do not edit by hand.
\begin{description}
  \item[\texttt{item\_popularity\_count}] How often an item was viewed in training sessions up to the cut-off. Drives the popularity baseline and the cold-start fallback ordering for items with no neighbour-session support.

  \item[\texttt{item\_category}] Catalogue attribute from the item-properties changelog. Known before any interaction, so it is available for items that never appear in training and can be used to break ties among equally popular candidates.

  \item[\texttt{visitor\_prior\_session\_count}] Number of training sessions a visitorid has started before the current one. This is an optional, strictly ablation-only signal: visitorid is a cookie, not a durable account, so the primary session-kNN method never depends on it, but a returning cookie is still weak positive evidence.

  \item[\texttt{session\_length\_so\_far}] Number of events already observed in the current session, i.e. the 0-indexed position itself. Short sessions get little context, which is exactly the "cold session" regime broken out in the evaluation report.

  \item[\texttt{time\_since\_session\_start}] Elapsed time between the session's first event and the *last observed* event (not the event being predicted). Captures browsing pace -- a visitor who has lingered for ten minutes behaves differently from one three seconds into a session.

  \item[\texttt{hour\_of\_day}] Hour extracted from the event being scored. Not an aggregate over history, so it carries no leakage risk: the wall-clock time of the prediction request is known by definition.

  \item[\texttt{is\_weekend}] Weekday/weekend split of browsing behaviour, used as an ablation feature alongside hour\_of\_day to test whether coarse temporal context improves candidate re-ranking over neighbour-vote scores alone.

\end{description}

